\documentclass[journal]{IEEEtran}

\usepackage{amsmath,amssymb}
\usepackage{booktabs}
\usepackage{array}
\usepackage{cite}
\usepackage{graphicx}
\usepackage{multirow}
\usepackage{threeparttable}
\usepackage{url}
\usepackage{xcolor}

% Keep full-width figures and tables visually attached to the surrounding text.
\setlength{\textfloatsep}{8pt plus 2pt minus 2pt}
\setlength{\dbltextfloatsep}{6pt plus 1pt minus 1pt}
\setlength{\floatsep}{6pt plus 2pt minus 2pt}
\setlength{\dblfloatsep}{6pt plus 2pt minus 2pt}

\newcommand{\auroc}{\mathrm{AUROC}}
\newcommand{\auprc}{\mathrm{AUPRC}}
\newcommand{\brier}{\mathrm{Brier}}
\newcommand{\ece}{\mathrm{ECE}}
\newcommand{\NA}{\textemdash}
\newcolumntype{L}[1]{>{\raggedright\arraybackslash}p{#1}}

\title{Multimodal COVID-19 Respiratory Sound Models: Temporal Stability and External Cough Transfer}

\author{Nishant Harkut and Ritu Tiwari
\thanks{The authors are with Atal Bihari Vajpayee Indian Institute of Information Technology and Management, Gwalior, India. Correspondence should be addressed to Nishant Harkut.}
}

\begin{document}
\bstctlcite{BSTcontrol}

\maketitle

\begin{abstract}
Respiratory-audio models can discriminate COVID-19 labels within development datasets, but transport across populations, collection processes, and calendar periods remains uncertain. We conducted a retrospective reliability study using Coswara as a multimodal source and COUGHVID as an independent cough-only target with semi-supervised status labels. The source workflow combined acoustic summaries with two standardized openSMILE paralinguistic representations, selected 800 source-ranked features, and evaluated four classifier families. The best observed cough--speech fusion achieved an area under the receiver operating characteristic curve (AUROC) of 0.897 in 314 participant-disjoint Coswara test participants. In modality-matched cough transfer, source-test AUROC values of 0.849--0.868 fell to 0.523--0.543 in 8,331 COUGHVID recording UUIDs. Declines were 0.310--0.345, and all independent-bootstrap 95\% confidence intervals had lower bounds of at least 0.247. Supporting self-supervised and spectrogram neural models yielded external AUROCs of 0.484 and 0.548. At 0.90 operational-label sensitivity, specificity was 0.112--0.161 and precision was 0.0347--0.0367, close to target-label prevalence 0.0342. Pseudo-label-supervised recalibration reduced probability error on held-out target recordings but did not recover discrimination. Secondary analyses found temporal instability, low overlap between early and late feature rankings, and strong source-label predictability from symptoms and collection context. Favorable source discrimination therefore did not transport to the evaluated external endpoint. These findings support frozen, modality-matched external evaluation with prevalence-aware operating points, calibration assessment, and explicit analysis of collection context.
\end{abstract}

\begin{IEEEkeywords}
COVID-19, dataset shift, external validation, respiratory audio, temporal stability.
\end{IEEEkeywords}

\section{Introduction}

\IEEEPARstart{R}{espiratory} sounds can be collected with ordinary microphones, which motivated their study as a low-cost screening signal during the COVID-19 pandemic. Public datasets and challenges rapidly enabled models based on cough, breath, speech, engineered acoustic descriptors, convolutional networks, and transformers \cite{laguarta2020cough,brown2020crowdsourced,muguli2021dicova,schuller2021compare,aytekin2024hst}. Several studies reported high discrimination under their chosen development protocols. Those results established technical feasibility, but they did not by themselves establish that a model would retain discrimination when recruitment, disease prevalence, devices, label construction, or recording instructions changed.

The distinction matters because crowdsourced biomedical data are not random samples from a stable target population. Infection status, symptoms, access to testing, hospital recruitment, geographic location, and calendar time can jointly affect participation. Audio quality and recording behavior may also vary with these factors. A model can therefore perform well by exploiting associations that are valid in one collection process but unstable elsewhere. This is a general problem in health artificial intelligence (AI), where dataset shift and acquisition-dependent shortcuts can cause an apparently accurate model to fail after deployment \cite{subbaswamy2020shift,quinonero2009datasetshift,ly2024shortcut,degrave2021shortcuts}.

Evidence specific to COVID-19 audio has already challenged optimistic benchmark interpretations. Han \emph{et al.} obtained an AUROC of 0.71 under participant-independent, demographically controlled evaluation, while less appropriate splits produced estimates as high as 0.90 \cite{han2022sounds}. In a larger polymerase-chain-reaction-referenced cohort, Coppock \emph{et al.} reported AUROC 0.846 before confounder matching and 0.619 after matching, and found no clear practical advantage over symptom-based prediction \cite{coppock2024symptoms}. Temporal drift has also been documented in dynamic cough data \cite{ganitidis2025drift}. These studies show that evaluation design is part of the scientific question, not an implementation detail.

We therefore treat model reliability, rather than model novelty, as the object of study. The primary question is whether a fixed cough representation and fixed model families trained on Coswara retain discrimination on an independently collected COUGHVID endpoint. Supporting questions ask whether the result persists across learned representations, whether probability recalibration can explain the external loss, whether source findings are stable over calendar time, and how much label information is present in symptoms and collection context. The study makes four contributions:

\begin{enumerate}
    \item A model- and modality-matched source-to-target transfer experiment across four conventional model families, with analysis-unit-level uncertainty for the change in discrimination.
    \item Representation-family checks using WavLM and CNN--BiGRU models evaluated on the same external cohort.
    \item An endpoint-reliability analysis that separates discrimination from calibration and reports prevalence-aware operating-point agreement.
    \item Secondary temporal, feature-stability, metadata, support-overlap, and incremental-value analyses that delimit plausible interpretations of the source-domain result.
\end{enumerate}

The article follows contemporary guidance for prediction-model reporting and external validation \cite{collins2024tripodai,moons2019probast,riley2024external1,riley2024external2}. It does not claim that the evaluated models are diagnostic devices, that the two datasets use identical reference standards, or that the cough-only external result validates multimodal fusion.

\section{Related Work}

\subsection{Respiratory-Audio Classification}

Early COVID-19 audio studies explored forced coughs, breathing, and speech with convolutional and transfer-learning systems \cite{laguarta2020cough,coppock2021pilot,pinkas2020voice,andreu2021generic}. Public resources broadened this work. Coswara provides multiple variants of cough, breathing, sustained vowels, and counting speech together with symptoms and demographics \cite{bhattacharya2023coswara}. COUGHVID provides a much larger cough-only collection with crowdsourced metadata and automated or expert annotations \cite{orlandic2021coughvid}. A subsequent semi-supervised study released consensus-based labels for a subset of that collection \cite{orlandic2023semisupervised}. DiCOVA and ComParE supplied common tasks and baselines \cite{muguli2021dicova,schuller2021compare,sharma2022dicova}.

Subsequent systems used ensembles, pretrained speech models, and spectrogram transformers \cite{mohammed2021ensemble,ponomarchuk2022achoo,aytekin2024hst,wang2022wavlm}. HST, for example, models local-to-global spectrogram context and reported strong within-dataset results on several respiratory-sound corpora \cite{aytekin2024hst}. Recent work has also compared deep neural trees and forests over multiple cough datasets \cite{islam2026dndf}. These studies address model construction or broad benchmark performance. Our question is different: we freeze a source workflow and quantify how its estimates change under an independent target collection.

\subsection{Realistic Validation and Clinical Interpretation}

Participant-independent splitting prevents recordings from one participant appearing in both development and test sets, but it does not remove recruitment or calendar associations. Han \emph{et al.} directly demonstrated sensitivity to split design and demographic bias \cite{han2022sounds}. Coppock \emph{et al.} used causal enrollment diagrams, matching, symptom comparators, and practical utility analyses to show how recruitment can inflate apparent audio performance \cite{coppock2024symptoms}. The present study extends this reliability perspective to a fixed Coswara-to-COUGHVID transfer and to calibration, operating points, feature stability, and source-context controls.

Reliability interpretation requires more than AUROC. Precision-recall performance depends on prevalence and is informative in imbalanced cohorts \cite{saito2015auprc}. Calibration asks whether predicted probabilities agree with observed frequencies \cite{vancalster2019calibration}. We therefore report discrimination, probability error, calibration error, and threshold-dependent quantities separately. Clinical utility would additionally require a clinically adjudicated outcome and a prospective decision context, neither of which is supplied by the external pseudo-label endpoint.

\section{Materials and Methods}

\subsection{Study Design and Analysis Hierarchy}

This was a retrospective secondary analysis of two de-identified, publicly documented respiratory-audio resources. Coswara was the source dataset for model development and source-domain evaluation. For the primary frozen-transfer analysis, all preprocessing choices, selected features, fitted models, and source thresholds were fixed before COUGHVID scoring. Separate descriptive operating-point and target-supervised recalibration analyses subsequently used COUGHVID labels.

We defined the model- and modality-matched cough transfer as the primary analysis. Source-domain multimodal fusion, learned-representation transfer, and target-label recalibration were supporting analyses. Time-stratified evaluation, retrospective early-to-late stress testing, feature stability, and source-target separability were secondary analyses. Metadata prediction, permutation controls, subgroup summaries, descriptive target operating points, and participant-aligned incremental-value comparisons were exploratory. The study was not preregistered. Source fusion rows are therefore described as best observed source results rather than confirmatory model comparisons. This reporting hierarchy distinguishes the central transfer question from post-hoc diagnostic analyses.

Figure~\ref{fig:workflow} separates source fitting from source robustness analyses and independent target scoring. COUGHVID entered only the external branch after source preprocessing choices, selected cough features, fitted classifiers, and source thresholds had been frozen.

\begin{figure*}[t]
    \centering
    \includegraphics[width=\textwidth]{figures/fig1_study_design.pdf}
    \caption{Study design and information flow. Coswara supplied multimodal source recordings for deterministic preprocessing, acoustic representation, source-only feature ranking, classifier fitting, participant aggregation, and validation-derived fusion. The fitted source workflow was evaluated under participant-disjoint, time-stratified, and retrospective early-to-late protocols. COUGHVID entered only the independent cough-transfer branch, where source-selected cough features, fitted models, and thresholds remained frozen. Reliability analyses were computed from the resulting predictions and audit tables.}
    \label{fig:workflow}
\end{figure*}

\subsection{Datasets, Labels, and Evaluation Units}

Coswara was collected from April 2020 to February 2022, primarily in India, and contains nine recordings per complete contribution: two coughs, two breaths, three sustained vowels, and two counting tasks \cite{bhattacharya2023coswara}. The analyzed snapshot contained 2,114 participants with a known operational binary label: 1,433 negative and 681 positive. Positive labels combined the available positive asymptomatic, mild, and moderate categories. Healthy records formed the negative class. Recovered, exposed, respiratory-illness, under-validation, and other unresolved categories were excluded from supervised endpoints. Labels were self-reported in the source data, although the dataset creators report that most positive participants were recruited through hospitals and health centers \cite{bhattacharya2023coswara}.

COUGHVID contains crowdsourced cough recordings, geographic and demographic fields, and status labels with heterogeneous provenance \cite{orlandic2021coughvid}. After the predetermined external inclusion and deduplication procedure, the target contained 8,331 unique recording UUIDs: 8,046 negative-labeled and 285 positive-labeled recordings, giving prevalence 0.0342. The public metadata provide one row per recording and do not supply a persistent person identifier, so repeated submissions by one individual cannot be clustered. The analyzed endpoint used the semi-supervised metadata field \texttt{status\_SSL}, which was produced by combining user labels with predictions from models trained on sparse expert annotations \cite{orlandic2023semisupervised}. Values designated COVID-19 or positive were mapped to the positive class, and values designated healthy or negative to the negative class. Because source-specific label mechanisms differ and the target label is not an RT-PCR reference standard, the transfer result is interpreted as recording-level transport and label agreement between operational dataset endpoints, not clinical diagnostic validation.

For the engineered models and WavLM, Coswara recording or segment probabilities were aggregated to participants. COUGHVID scores were evaluated per recording UUID because no persistent person identifier was available. The supporting CNN--BiGRU check used recording-level scores at both source and target. Table~\ref{tab:cohorts} distinguishes dataset role, modality, label source, and evaluation unit.

\begin{table*}[t]
\caption{Cohorts and Evaluation Roles}
\label{tab:cohorts}
\centering
\footnotesize
\begin{threeparttable}
\begin{tabular}{L{2.1cm}L{2.2cm}L{2.4cm}L{2.2cm}L{2.1cm}L{2.7cm}}
\toprule
Dataset or protocol & Role & Available input & Analyzed units & Positive prevalence & Interpretation boundary \\
\midrule
Coswara known-label snapshot & Source development and internal evaluation & Cough, breath, sustained vowels, counting speech & 2,114 (681 positive) & 0.322 & Crowdsourced operational labels. Unresolved and recovered categories excluded \\
COUGHVID external cohort & Independent target & Cough only & 8,331 recording UUIDs (285 positive) & 0.034 & Different collection and label mechanism. No persistent person identifier or external breath or speech \\
Existing participant split & Source-domain support & Participant-disjoint train, validation, test across collection dates & 1,479 / 317 / 318 & 0.322 / 0.322 / 0.324 & Source discrimination, not evidence of external transport \\
Time-stratified participant split & Source robustness & Participant-disjoint sets stratified across collection time & 1,258 / 412 / 444 & 0.319 / 0.333 / 0.322 & Different complete-case counts by modality and fusion \\
Early-to-late stress test & Secondary temporal analysis & Train through 2020-12-04, validate through 2021-07-29, and test later recordings & 1,268 / 422 / 424 & 0.079 / 0.569 / 0.804 & Retrospective and reuses source-selected representation. Not prospective validation \\
\bottomrule
\end{tabular}
\begin{tablenotes}
\item Participant counts in multimodal analyses can be lower because a participant must have predictions for every fused modality.
\end{tablenotes}
\end{threeparttable}
\end{table*}

\subsection{Audio Quality Control and Preprocessing}

Audio was decoded as mono floating-point waveforms and resampled to 16~kHz. Files with unreadable content, insufficient duration, excessive silence, or clipping were flagged before feature extraction. The audit classified 23,539 recordings as usable, 597 as mostly silent, 390 as corrupt, and 186 as short. Silence trimming, energy-based event detection, peak normalization, and fixed-duration cropping or padding were deterministic. Event timing and active-audio ratio were retained as acquisition-context descriptors rather than interpreted as disease-specific physiology. Full settings are reported in Supplementary Table~II.

\subsection{Acoustic Representations and Feature Selection}

The engineered table combined three representations. First, a broad acoustic bank summarized 40 mel-frequency cepstral coefficients (MFCCs), first and second MFCC differences, 64 log-mel bands, chroma, spectral contrast, tonal centroid features, energy, zero-crossing rate, spectral centroid, bandwidth, roll-off, flatness, and a tempogram using librosa-based extraction \cite{mcfee2015librosa}. For each trajectory, distributional summaries captured location, spread, quantiles, extrema, and temporal variation. Second and third, openSMILE produced the ComParE 2016 and IS10 paralinguistic sets \cite{eyben2010opensmile,schuller2010is10,schuller2016compare}. After identifier and context columns, the merged table contained 10,147 columns before selection.

Nonfinite feature values were replaced by zero before ranking. On source-training rows, a LightGBM ranker produced gain importances separately for each modality. Let $g_{jm}$ denote the gain assigned to feature $j$ in modality $m$. The cross-modality ranking score was
\begin{equation}
    q_j=\frac{1}{M}\sum_{m=1}^{M}
    \frac{g_{jm}}{s_m}, \qquad
    \mathcal{S}_{800}=\operatorname{Top}_{800}\{q_j\}.
\end{equation}
Here, $s_m=\max_{\ell}g_{\ell m}$ when the maximum gain is positive and $s_m=1$ otherwise. Only modalities containing both classes contributed to the average. Normalization prevented a modality with a larger raw gain scale from dominating the shared ranking. The external cohort used the frozen set $\mathcal{S}_{800}$ with no target-driven reselection. The retrospective temporal stress test also reused this list rather than reselecting solely within its early training period. Within each conventional classifier pipeline, a variance filter removed constant columns and a univariate ANOVA filter was fitted on the training partition. The tree-boosting pipelines retained the top 80\% of the selected features, while RBF-SVC retained 60\% after standardization.

Two learned representations were included as supporting architecture checks. WavLM-base-plus, a self-supervised speech transformer \cite{wang2022wavlm}, processed up to four overlapping 3-s cough segments per recording and was fine-tuned on source data. Its segment scores were averaged per Coswara participant and per COUGHVID recording UUID. A CNN--BiGRU model processed one log-mel spectrogram per recording with convolutional filters followed by a bidirectional gated recurrent unit. Its metrics remained recording-level. These branches used their own training recipes and analysis units, so their absolute source scores were not pooled with the engineered-model estimates.

\subsection{Model Families, Imbalance Handling, and Fusion}

The controlled engineered-feature experiment used LightGBM, XGBoost, CatBoost, and calibrated RBF-SVC \cite{ke2017lightgbm,chen2016xgboost,prokhorenkova2018catboost}. Tree-boosting models used synthetic minority oversampling with three neighbors on training data only \cite{chawla2002smote}. RBF-SVC used class balancing, standard scaling, and three-fold sigmoid calibration \cite{platt1999probabilistic,niculescu2005calibration}. All model settings were fixed before external scoring and are listed in Supplementary Table~V.

Within a modality, recording-level probabilities were averaged for participant $i$:
\begin{equation}
    p_{im}=\frac{1}{R_{im}}\sum_{r=1}^{R_{im}}p_{irm},
\end{equation}
where $m$ denotes modality and $R_{im}$ is the number of available recordings. Multimodal source experiments considered three validation-only fusion rules. Uniform fusion averaged modality probabilities. Validation-weighted fusion used
\begin{equation}
    \tilde{w}_{m}=\max(\auprc_{m}^{(val)}-0.5,0.01), \qquad
    w_m=\frac{\tilde{w}_m}{\sum_j\tilde{w}_j},
\end{equation}
and $p_i=\sum_m w_m p_{im}$. Stacked fusion used $p_i^{(stack)}=\sigma(\beta_0+\sum_m\beta_m p_{im})$, where $\sigma$ is the logistic link and the coefficients were fitted by class-balanced logistic regression on validation participants with all required modalities. Neither averaging weights nor stacking coefficients were chosen from the source test set. For each candidate, the operating threshold maximized the mean of validation sensitivity and specificity. Ties were resolved toward the threshold closest to 0.5, and the result was fixed for source test evaluation.

\subsection{Evaluation Protocols}

\subsubsection{Source-domain participant split}
Participants were assigned to 70\% training, 15\% validation, and 15\% test groups with stratification by the binary label. All recordings from one participant remained in one group. This protocol estimated source-domain discrimination and supported model fitting, validation-derived fusion, and threshold setting.

\subsubsection{Time-stratified participant split}
A second participant-disjoint protocol distributed train, validation, and test participants across the collection timeline while preserving similar prevalence. It tested sensitivity to a different source split, not future deployment.

\subsubsection{Retrospective early-to-late stress test}
Training included recordings through 4 December 2020. Validation covered the subsequent period through 29 July 2021, and testing covered 4 August 2021 through 24 February 2022. This analysis deliberately stressed calendar stability. It reused the frozen source representation selected in the broader source workflow. It is therefore described as a retrospective stress test, not leakage-free prospective validation.

\subsubsection{External transfer}
Models trained on Coswara cough recordings were applied without refitting to the COUGHVID cohort. The engineered comparison held the selected feature set, preprocessing, model family, and source validation threshold fixed. WavLM and CNN--BiGRU were scored as separate learned-representation checks. COUGHVID was never mixed into source training.

\subsection{Outcomes, Uncertainty, and Reliability Analyses}

The principal discrimination measures were AUROC and AUPRC. We also calculated balanced accuracy, sensitivity, specificity, F1 score, negative log likelihood, Brier score \cite{brier1950verification}, and expected calibration error (ECE) \cite{vancalster2019calibration}. For labels $y_i\in\{0,1\}$ and probabilities $p_i$,
\begin{equation}
\begin{split}
    \brier&=\frac{1}{n}\sum_{i=1}^{n}(p_i-y_i)^2,\\
    \ece&=\sum_{b=1}^{10}\frac{n_b}{n}\left|\bar{p}_b-\bar{y}_b\right|,
\end{split}
\end{equation}
where $n_b$, $\bar{p}_b$, and $\bar{y}_b$ are the size, mean prediction, and observed positive fraction in the $b$th equal-width probability bin. AUROC evaluates ranking, whereas Brier score and ECE evaluate the probability scale. Because AUPRC and precision depend on prevalence, every external interpretation used the observed target positive-label prevalence as a reference \cite{saito2015auprc}.

For matched source-to-target model families, Coswara participants and COUGHVID recording UUIDs are independent across datasets. We therefore obtained the confidence interval for
\begin{equation}
    \Delta=\auroc_{source}-\auroc_{target}
\end{equation}
by independently resampling source participants and target recording UUIDs with replacement and recomputing the difference. We used 2,000 replicates, a base random seed of 42, and the 2.5th and 97.5th percentiles of the bootstrap distribution. This is an independent two-sample bootstrap, not subtraction of two marginal confidence intervals \cite{carpenter2000bootstrap}. DeLong tests were reserved for model comparisons evaluated on the same participants \cite{delong1988roc}.

On the external score distribution, descriptive operating-point analysis chose the threshold with highest specificity among those attaining sensitivity of at least 0.90 against \texttt{status\_SSL}. This point quantifies the false-positive burden required for high operational-label sensitivity. It is not a clinical sensitivity estimate or a deployable target-tuned threshold. A separate recalibration experiment split COUGHVID into 2,082 calibration UUIDs and 6,249 evaluation UUIDs. Platt and isotonic mappings were fitted only on the calibration subset, then evaluated on the held-out subset \cite{platt1999probabilistic,niculescu2005calibration}. Platt scaling used logistic regression on source-model log-odds, whereas isotonic regression used raw probabilities and could introduce ties. The experiment tests whether probability-scale correction recovers discrimination against the operational target label rather than assuming AUROC must be exactly invariant.

Secondary analyses included repeated-seed training, a source-versus-target domain classifier, and overlap of two independently ranked top-800 feature sets from early and late periods. The latter used LightGBM gain rankings and Jaccard similarity $|A\cap B|/|A\cup B|$. The domain classifier was a standardized, class-balanced logistic regression evaluated with five-fold out-of-fold predictions on the 500 common features with highest source-training variance. Exploratory context models used standardized, class-balanced logistic regression to predict the source label from symptoms, demographics, date, country, duration, sample rate, and quality fields. Models were retrained after label permutation as an implementation-leakage control. Grouped permutation importance quantified the reduction in predictive performance after shuffling each variable group \cite{fisher2019modelreliance}. Participant-aligned logistic models compared metadata-only, audio-only, and combined predictions on complete cases. Effect sizes, bootstrap intervals, and paired DeLong tests were reported. These small complete-case comparisons remained exploratory.

\section{Results}

\subsection{Cohorts and Source-Domain Performance}

The source and target differed substantially in class composition (Table~\ref{tab:cohorts}). Positive prevalence was 0.322 among known-label Coswara participants and 0.034 among external COUGHVID recordings. In the source-domain participant split, the best observed cough--speech stacked logistic fusion achieved AUROC 0.897, AUPRC 0.863, balanced accuracy 0.825, and F1 score 0.752 on 314 complete-case test participants. Across four seeds for this fusion, mean AUROC was $0.895\pm0.003$. This indicates low AUROC variation across those retraining seeds. It does not quantify cohort-sampling uncertainty. The estimate is not directly comparable with the cough-only target because COUGHVID has no breath or speech recordings.

The best observed time-stratified source row was a three-modality uniform mean with AUROC 0.849, AUPRC 0.783, and ECE 0.045 on 431 complete-case participants. Its model and modality composition differed from the source headline. Consequently, these values are reported as separate source robustness estimates rather than consecutive steps in one degradation curve.

\begin{figure*}[t]
    \centering
    \includegraphics[width=\textwidth]{figures/fig2_external_transport.pdf}
    \caption{Source-to-target transport. (a) AUROC at each dataset's available analysis-unit level for four engineered-feature models and two learned-representation checks. Lines connect the same model family, not identical training recipes across families. The gray dashed line marks chance AUROC. (b) AUROC difference for the controlled engineered cough comparison. Confidence intervals independently resample source participants and target recording UUIDs. (c) External AUPRC. The orange dashed line marks COUGHVID prevalence (0.0342).}
    \label{fig:transport}
\end{figure*}

\subsection{Controlled Cough-Only External Transfer}

Every engineered model family showed strong source-test discrimination and near-chance external discrimination (Table~\ref{tab:external}, Fig.~\ref{fig:transport}(a)). Source-test AUROC ranged from 0.849 to 0.868. External AUROC ranged from 0.523 to 0.543. The corresponding AUPRC values fell from 0.780--0.812 to 0.037--0.040, close to the target prevalence.

The independently bootstrapped AUROC decline was 0.310 for LightGBM, 0.318 for CatBoost and XGBoost, and 0.345 for RBF-SVC. All 95\% confidence intervals excluded zero, with lower bounds between 0.247 and 0.291 (Fig.~\ref{fig:transport}(b)). The direction and magnitude of decline were similar across all four conventional classifier families.

\begin{table*}[t]
\caption{Controlled Cough-Only Source-to-Target Transfer}
\label{tab:external}
\centering
\footnotesize
\begin{threeparttable}
\begin{tabular}{lccccccc}
\toprule
Model & $n_{source}$ & Source AUROC & Source AUPRC & $n_{target}$ & Target AUROC & Target AUPRC & AUROC decline (95\% CI) \\
\midrule
CatBoost & 316 & 0.849 & 0.788 & 8,331 & 0.531 & 0.040 & 0.318 (0.256--0.376) \\
LightGBM & 316 & 0.853 & 0.797 & 8,331 & 0.543 & 0.040 & 0.310 (0.247--0.366) \\
RBF-SVC & 316 & 0.868 & 0.812 & 8,331 & 0.523 & 0.037 & 0.345 (0.291--0.398) \\
XGBoost & 316 & 0.850 & 0.780 & 8,331 & 0.532 & 0.039 & 0.318 (0.262--0.371) \\
\midrule
WavLM & 316 & 0.812 & 0.734 & 8,331 & 0.484 & 0.032 & \NA \\
CNN--BiGRU & 636 & 0.737 & 0.574 & 8,331 & 0.548 & 0.044 & \NA \\
\bottomrule
\end{tabular}
\begin{tablenotes}
\item The first four rows use participant-level source scores and recording-UUID-level target scores with the same frozen 800-feature cough representation and model family. WavLM uses the same source and target analysis units. CNN--BiGRU uses 636 source recordings and 8,331 target recordings. The learned branches have different training recipes, and no bootstrap interval was computed for their decline.
\end{tablenotes}
\end{threeparttable}
\end{table*}

The two supporting learned-representation analyses also produced external AUROC point estimates near 0.5. WavLM declined from source-test AUROC 0.812 to external AUROC 0.484. CNN--BiGRU declined from 0.737 to 0.548. Their external AUPRC values were 0.032 and 0.044. Because uncertainty intervals were not computed and training recipes and source denominators differed, these are qualitative consistency checks rather than an architecture ranking.

\begin{figure*}[t]
    \centering
    \includegraphics[width=\textwidth]{figures/fig3_endpoint_reliability.pdf}
    \caption{External operational-label diagnostics. (a) Specificity against \texttt{status\_SSL} at a descriptively chosen threshold with sensitivity at least 0.90. (b) Precision against the same pseudo-label endpoint. The dashed line is positive-label prevalence. Both operating-point axes begin at zero. (c) Held-out COUGHVID evaluation before and after pseudo-label-supervised Platt or isotonic recalibration. AUROC, Brier score, and ECE are displayed as separate metric rows on one linear scale. Platt scaling reduces probability error without changing AUROC. These panels characterize agreement with the operational label, not clinical diagnostic utility.}
    \label{fig:labelagreement}
\end{figure*}

\subsection{External Label Agreement and Recalibration}

At a threshold chosen to attain sensitivity of at least 0.90 against \texttt{status\_SSL}, specificity ranged from 0.112 to 0.161 across the four controlled model families (Fig.~\ref{fig:labelagreement}(a)). Precision ranged from 0.0347 to 0.0367, only 0.0005--0.0025 above the observed positive-label prevalence of 0.0342 (Fig.~\ref{fig:labelagreement}(b)). Thus, high agreement sensitivity required labeling most negative-labeled recordings as positive and produced little enrichment above the target pseudo-label base rate. These quantities do not estimate clinical screening utility.

On the held-out external recalibration experiment, the unrecalibrated LightGBM model had AUROC 0.539, Brier score 0.0756, and ECE 0.0992. Platt scaling reduced Brier score to 0.0330 and ECE to 0.00017, while AUROC remained 0.539. Isotonic regression produced AUROC 0.532, Brier score 0.0331, and ECE 0.00027. Recalibration therefore reduced measured probability error but did not improve ranking discrimination against the operational label (Fig.~\ref{fig:labelagreement}(c)).

\begin{figure*}[t]
    \centering
    \includegraphics[width=\textwidth]{figures/fig4_context_and_stability.pdf}
    \caption{Secondary context and stability analyses. (a) Source label prediction from metadata configurations and retraining after label permutation. (b) Grouped permutation importance for the full symptoms-plus-context model. (c) Overlap of independently ranked early and late top-800 feature sets. (d) Positive-label prevalence across the retrospective early-to-late cohorts. These panels are exploratory diagnostics and should not be interpreted as causal attribution to the audio model or biological disease progression.}
    \label{fig:context}
\end{figure*}

\subsection{Temporal Stability and Collection Context}

The retrospective early-to-late analysis selected a breath ensemble and achieved AUROC 0.698, Brier score 0.650, and ECE 0.711 on 411 complete-case test participants. Across three seeds, mean AUROC was $0.691\pm0.007$. Interpretation is constrained by the prevalence change from 0.079 in training to 0.569 in validation and 0.804 in test (Table~\ref{tab:cohorts}, Fig.~\ref{fig:context}(d)), and by reuse of the broader source-selected representation. The high temporal-test AUPRC of 0.896 must therefore be viewed against the test prevalence of 0.804 rather than compared directly with lower-prevalence endpoints.

In a separate exploratory split at 13 August 2020, early and late LightGBM rankings shared 110 of their top 800 features. The union contained 1,490 features, giving Jaccard similarity 0.074 (Fig.~\ref{fig:context}(c)). The early and late samples also differed sharply in positive count. This result identifies feature-ranking instability under combined calendar, prevalence, and recruitment change but does not isolate a biological drift mechanism.

Source labels were highly predictable from symptoms and collection context. A full context model reached AUROC 0.964. After retraining on 20 independently permuted label vectors, mean AUROC was 0.503, with permutation range 0.465--0.559 (Fig.~\ref{fig:context}(a)). Symptoms-only AUROC was 0.932, while demographic and protocol fields produced 0.914. Grouped permutation importance in the full context configuration assigned 71.7\% of total importance to recording-protocol variables, 18.5\% to symptoms, 5.2\% to comorbidity, and 4.6\% to demographics. Recording year was the leading individual variable (Fig.~\ref{fig:context}(b)). These are predictive associations in the dataset and do not prove that a specific audio model encoded the same variables.

The source-versus-target domain classifier achieved AUROC 0.750. Its source probability band, defined by the fifth and 95th percentiles of source-domain scores, excluded 25.2\% of external recording-level observations. This probability-band diagnostic supports covariate shift but is not a formal proof of a causal positivity violation.

\subsection{Incremental Value of Audio on Aligned Participants}

Complete-case alignment between metadata and final audio predictions was limited to 58--61 test participants. On 61 participants, symptoms alone produced AUROC 0.888, the best aligned audio candidate produced 0.818, and the combined model produced 0.951. The AUROC increment over symptoms was 0.063 (95\% bootstrap CI $-0.005$ to 0.149, paired DeLong $p=0.104$). Full context metadata produced AUROC 0.924, while full context plus the selected audio fusion produced 0.933, an increment of 0.009 (95\% CI $-0.026$ to 0.040, $p=0.538$). These small-sample estimates are inconclusive for symptoms and show little measured increment beyond the full context model (Table~\ref{tab:support}).

\begin{table*}[t]
\caption{Supporting and Exploratory Reliability Results}
\label{tab:support}
\centering
\footnotesize
\begin{tabular}{L{3.2cm}L{2.5cm}L{2.0cm}L{2.0cm}L{4.1cm}}
\toprule
Analysis & Endpoint & Sample size & Estimate & Boundary on interpretation \\
\midrule
Best observed source multimodal fusion & AUROC / AUPRC & 314 & 0.897 / 0.863 & Cough--speech source result. Not externally validated as multimodal \\
Source fusion across seeds & Mean AUROC $\pm$ SD & 4 seeds & $0.895\pm0.003$ & Training variability, not cohort-sampling uncertainty \\
Time-stratified source fusion & AUROC / ECE & 431 & 0.849 / 0.045 & Different three-modality fusion and complete-case cohort \\
Early-to-late stress test & AUROC / ECE & 411 & 0.698 / 0.711 & Retrospective with large prevalence change and a frozen broader source representation \\
Early/late feature overlap & Jaccard similarity & 800 features per period & 0.074 & Separate exploratory date split. Prevalence and recruitment differ \\
Full context metadata & Observed / shuffled AUROC & 20 permutations & 0.964 / 0.503 & Predictive association, not proof of audio-model shortcut use \\
Symptoms plus audio & AUROC increment & 61 & 0.063 ($-0.005$, 0.149) & Exploratory complete-case comparison. DeLong $p=0.104$ \\
Full context plus audio & AUROC increment & 61 & 0.009 ($-0.026$, 0.040) & Exploratory complete-case comparison. DeLong $p=0.538$ \\
\bottomrule
\end{tabular}
\end{table*}

\section{Discussion}

\subsection{Principal Findings}

This study found a large and consistent gap between source-domain and external discrimination. When modality, engineered representation, and model family were held fixed, four cough models declined by 0.310--0.345 AUROC from Coswara to COUGHVID. Every confidence interval remained far from zero, and external AUPRC remained near the 3.42\% target prevalence. WavLM and CNN--BiGRU provided qualitatively consistent external point estimates under their respective source-trained recipes. The controlled result therefore is not a multimodal-to-cough comparison and was reproduced across four conventional classifier families.

The operating-point and recalibration results refine that conclusion. At sensitivity at least 0.90 against the pseudo-label endpoint, specificity was at most 0.161 and precision was barely above positive-label prevalence. Pseudo-label-supervised calibration substantially improved probability error but did not improve AUROC. The dominant external problem was thus poor ranking discrimination against the evaluated operational endpoint, not only an incorrect threshold or probability scale. These analyses cannot establish clinical utility because the target was not diagnostically adjudicated.

The secondary analyses identify plausible sources of instability without claiming a single cause. The analyzed Coswara timeline combined a large prevalence change with low overlap between early and late feature rankings. Source labels were strongly associated with symptoms and collection context, especially recording year, and a domain classifier separated source observations from target recordings. These analyses document label--context association and measurable source--target separability. They do not establish which context variables, if any, were used by the audio models, nor do they establish that no portable acoustic information exists.

\subsection{Relation to Prior Evidence}

The results reinforce the realistic-evaluation findings of Han \emph{et al.}, who showed that sample overlap and demographic imbalance can raise apparent performance \cite{han2022sounds}. They also complement Coppock \emph{et al.}, who found a reduction from 0.846 before matching to 0.619 after measured-confounder matching and no practical improvement over symptoms \cite{coppock2024symptoms}. Our external experiment differs in using an independent public target and in holding cough modality and four model families fixed. The small participant-aligned incremental analysis points in the same cautious direction as Coppock \emph{et al.}, but its uncertainty is too wide to claim that symptoms and audio are equivalent.

Ganitidis \emph{et al.} explicitly modeled drift and adaptation in COVID-19 cough streams \cite{ganitidis2025drift}. We did not perform online adaptation. Instead, our retrospective temporal test documents the instability that would motivate monitoring, while its prevalence and feature-selection limitations prevent a prospective claim. More generally, the source-target separability and performance loss are consistent with health-AI work on acquisition-dependent shortcuts and dataset shift \cite{subbaswamy2020shift,ly2024shortcut}. Robust cough systems increasingly incorporate out-of-distribution detection \cite{chen2023ood}. Our findings support evaluating such safeguards against clinically defined external targets rather than only generic non-cough sounds.

\subsection{Implications for Respiratory-Audio Evaluation}

Three evaluation practices follow from these results. First, model comparisons should identify the prediction target and hold modality constant. The relevant external comparison here is cough-to-cough, not multimodal-to-cough. Second, independent target data should remain untouched until the source workflow is frozen. Recalibration or target adaptation can then be studied explicitly, with separate target calibration and evaluation subsets. Third, discrimination, calibration, prevalence, and operating-point behavior answer different questions. A high AUROC does not imply calibrated probabilities, and operating-point metrics against a pseudo-label do not establish clinical utility.

The study also shows why context variables require transparent handling. Duration, timing, date, device characteristics, and quality may be necessary for signal processing or may carry collection-specific information. Reporting their inclusion and testing label predictability from context allow readers to distinguish measured label--context association from direct evidence about audio-model behavior.

\subsection{Strengths}

The primary analysis uses an explicit source-to-target direction, a frozen representation, the same cough modality, four distinct classifier families, participant-level source aggregation, recording-level target evaluation, and independent-bootstrap intervals for the performance change. Learned-representation point estimates provide a qualitative check beyond engineered descriptors. Separate target calibration and evaluation subsets distinguish probability correction from discrimination against the operational label. Negative-control retraining after label permutation checks the metadata pipeline. Finally, the analysis hierarchy and evidence map separate primary results from post-hoc diagnostics.

\subsection{Limitations}

The study has important limitations. First, Coswara and COUGHVID do not provide identical clinical reference standards. Coswara includes self-reported test categories with substantial hospital-linked recruitment, while the COUGHVID endpoint uses the semi-supervised \texttt{status\_SSL} pseudo-label. The external result therefore measures dataset-endpoint transport and label agreement, not transport between two identically adjudicated diagnostic cohorts.

Second, COUGHVID is cough-only. It externally tests the cough branch and cannot validate the full cough--breath--speech fusion. The source multimodal result and external cough result are deliberately not subtracted. A multimodal external cohort with harmonized labels would be required to assess fusion transport.

Third, the temporal analysis is retrospective and reuses a representation frozen in the broader source workflow. Its cohorts also differ sharply in prevalence and recruitment period. The estimate is evidence of instability under a combined calendar stress, not a leakage-free prospective estimate or a causal effect of viral evolution. The early-versus-late feature overlap uses a separate date boundary and is exploratory.

Fourth, source and target differ simultaneously in geography, prevalence, collection interface, devices, instructions, and label construction. The study quantifies the joint transport failure but cannot assign proportions of that failure to individual mechanisms. The probability-band overlap diagnostic is not a formal support theorem.

Fifth, context models include symptoms, demographics, date, country, recording duration, sample rate, and quality. Their high AUROC demonstrates label-context association, not that the audio models encoded every listed field. Direct causal attribution would require a design that intervenes on or carefully matches those factors.

Sixth, COUGHVID identifies recordings by UUID but does not expose a persistent person identifier. The target analysis is therefore recording-level, and possible repeated submissions by one person could not be grouped in resampling or evaluation. Participant-aligned incremental-value analyses in Coswara had only 58--61 complete cases. Their intervals are wide, so they cannot resolve modest value beyond symptoms. Target recalibration used labeled COUGHVID data and should be interpreted as supervised site adaptation, not an unsupervised deployment procedure. No model was evaluated prospectively, no clinical workflow was implemented, and no claim of diagnostic readiness is made.

\section{Conclusion}

A respiratory-audio workflow that discriminated COVID-19 labels within Coswara did not retain useful discrimination on the evaluated COUGHVID pseudo-label endpoint. The AUROC decline was similar across four matched conventional model families, and two learned-representation checks produced qualitatively consistent external point estimates. Pseudo-label-supervised recalibration reduced calibration error but did not improve ranking. Temporal, feature-stability, metadata, and domain-separation analyses further show why source performance must be interpreted within its collection process. Reliable evidence for respiratory-audio screening therefore requires frozen source-to-target evaluation, harmonized clinical endpoints, prevalence-aware operating points, calibration assessment, and explicit analysis of recruitment and acquisition context.

\section*{Ethics Statement}

No new participants were recruited. This work is a secondary analysis of de-identified public research datasets. The original Coswara and COUGHVID publications describe their respective consent, ethics, and data-governance procedures \cite{bhattacharya2023coswara,orlandic2021coughvid}.
% Submission check: confirm whether a local secondary-analysis determination is required.

\section*{Data and Code Availability}

Coswara and COUGHVID are available under the terms specified by their original maintainers \cite{bhattacharya2023coswara,orlandic2021coughvid}. Analysis code, frozen configuration files, derived aggregate tables, and figure-generation scripts are maintained in the project repository. Participant-level derived predictions should be released only where permitted by the source dataset terms.

\section*{Author Contributions}

N.H. performed data curation, software implementation, formal analysis, visualization, and drafting. R.T. supervised the study and contributed to study design and interpretation.
% Submission check: all authors must approve the final CRediT wording.

\section*{Conflict of Interest}

The authors declare no conflicts of interest.
% Submission check: all authors must confirm the disclosure statement.

\bibliographystyle{IEEEtran}
\bibliography{references}

\end{document}
